\subsection{Biểu diễn không gian bằng Laplacian Positional Encoding}
\label{subsec:laplacian_pe}

Trong các mô hình Transformer tiêu chuẩn, cơ chế self-attention không tự thân nắm bắt được
thông tin vị trí hay cấu trúc của các phần tử đầu vào, mà cần được bổ sung thông qua các cơ
chế mã hóa vị trí (positional encoding). Đối với dữ liệu chuỗi, vị trí có thể được xác định
một cách tự nhiên theo trật tự thời gian. Tuy nhiên, trong bài toán giao thông trên đồ thị,
các tuyến đường không được sắp xếp theo một trật tự tuyến tính cố định, mà tồn tại trong
một không gian phi Euclid được xác định bởi cấu trúc mạng lưới giao thông.

Do đó, để cung cấp thông tin không gian cho mô hình Transformer, nhóm sử dụng
\textit{Laplacian Positional Encoding}, một phương pháp mã hóa vị trí dựa trên phổ của đồ
thị, đã được chứng minh là hiệu quả trong các mô hình Transformer cho dữ liệu đồ thị.

\subsubsection{Xây dựng đồ thị không gian trong vùng giao thông}
\label{subsubsec:zone_graph}

Đối với mỗi vùng giao thông chức năng được xây dựng ở các bước trước, nhóm xem vùng
này như một đồ thị con $G_z = (V_z, E_z)$ của đồ thị giao thông toàn cục, trong đó mỗi đỉnh
$v_i \in V_z$ tương ứng với một tuyến đường trong vùng, và các cạnh $E_z$ được xác định
dựa trên quan hệ kết nối topology giữa các tuyến trong mạng lưới giao thông.

Từ đồ thị con này, một ma trận kề $A_z \in \mathbb{R}^{|V_z| \times |V_z|}$ được xây dựng,
trong đó $A_z(i,j) = 1$ nếu tồn tại cạnh nối giữa hai tuyến $i$ và $j$ trong đồ thị giao
thông, và $A_z(i,j) = 0$ trong trường hợp ngược lại. Trên cơ sở đó, ma trận bậc
$D_z$ được xác định với $D_z(i,i) = \sum_j A_z(i,j)$.

\subsubsection{Laplacian và phân rã phổ}
\label{subsubsec:laplacian_eigendecomposition}

Ma trận Laplacian của đồ thị vùng được định nghĩa dưới dạng Laplacian tổ hợp:
\begin{equation}
L_z = D_z - A_z.
\end{equation}

Ma trận Laplacian này mã hóa đầy đủ cấu trúc kết nối của vùng giao thông và đóng vai trò
như một toán tử khuếch tán trên đồ thị. Thực hiện phân rã phổ của $L_z$ cho phép ta thu
được hệ eigenvalue và eigenvector:
\begin{equation}
L_z = U_z \Lambda_z U_z^\top,
\end{equation}
trong đó $U_z = [\mathbf{u}_1, \mathbf{u}_2, \dots, \mathbf{u}_{|V_z|}]$ là các eigenvector
tương ứng với các eigenvalue được sắp xếp theo thứ tự không giảm
$0 = \lambda_1 \le \lambda_2 \le \dots$.

Trong nghiên cứu này, nhóm chỉ sử dụng $d_{\text{spa}}$ eigenvector ứng với các
eigenvalue nhỏ nhất (bỏ qua eigenvector hằng ứng với $\lambda_1 = 0$ nếu cần), tạo thành
ma trận mã hóa vị trí không gian:
\begin{equation}
\mathbf{P}_z = [\mathbf{u}_2, \mathbf{u}_3, \dots, \mathbf{u}_{d_{\text{spa}}+1}]
\in \mathbb{R}^{|V_z| \times d_{\text{spa}}}.
\end{equation}

Việc lựa chọn các eigenvector bậc thấp tương ứng với các thành phần tần số thấp của đồ
thị cho phép mã hóa cấu trúc không gian một cách trơn (smooth) và mang tính toàn cục,
phản ánh vai trò cấu trúc của mỗi tuyến đường trong vùng giao thông.

\subsubsection{Ý nghĩa của eigenvector bậc thấp trong bối cảnh giao thông}
\label{subsubsec:low_freq_interpretation}

Các eigenvector bậc thấp của Laplacian đồ thị có thể được xem như các ``tọa độ cấu trúc''
(structural coordinates) của các tuyến đường trong vùng. Các tuyến có vị trí cấu trúc
gần nhau trong đồ thị (ví dụ cùng nằm ở trung tâm vùng hoặc cùng thuộc rìa vùng) sẽ có
biểu diễn gần nhau trong không gian phổ, ngay cả khi chúng không kề nhau trực tiếp về
mặt topology.

Trong bối cảnh giao thông đô thị, điều này đặc biệt quan trọng vì hiện tượng ùn tắc thường
lan truyền một cách liên tục và có cấu trúc theo vùng, thay vì thay đổi đột ngột trên từng
tuyến riêng lẻ. Việc sử dụng các eigenvector bậc thấp giúp mô hình nắm bắt được cấu trúc
không gian tổng thể của vùng giao thông, đồng thời giảm độ nhạy với nhiễu cục bộ hoặc các
liên kết không ổn định.

\subsubsection{Tích hợp Laplacian Positional Encoding vào Transformer}
\label{subsubsec:pe_integration}

Ma trận mã hóa vị trí không gian $\mathbf{P}_z$ được sử dụng như một \textit{spatial
positional encoding} cho các tuyến đường trong vùng giao thông. Đối với mỗi mẫu huấn
luyện, mã hóa này được ghép với biểu diễn theo thời gian của dữ liệu đầu vào và được cộng
vào embedding đầu vào của mô hình Transformer.

Bằng cách này, cơ chế self-attention trong Transformer không chỉ dựa trên sự tương đồng
ngữ nghĩa của các đặc trưng giao thông theo thời gian, mà còn được điều kiện hóa bởi thông
tin vị trí cấu trúc của các tuyến trong vùng. Điều này cho phép mô hình học được các tương
tác không gian--thời gian phức tạp và vượt qua các giới hạn của các mô hình dựa trên cơ chế
truyền thông tin cục bộ truyền thống.

Laplacian Positional Encoding do đó đóng vai trò then chốt trong việc cung cấp thông tin
không gian cho mô hình Transformer encoder-only, tạo nền tảng cho việc học biểu diễn
vùng giao thông một cách hiệu quả và có ý nghĩa về mặt cấu trúc.
